\documentclass[
12pt,
a4paper,
svgnames,
]{tufte-book}
\usepackage{fontspec}
\usepackage{polyglossia}
\setdefaultlanguage{thai}
\setcounter{secnumdepth}{2}
\setcounter{tocdepth}{1}

\XeTeXlinebreaklocale "th_TH"
\XeTeXlinebreakskip = 0pt plus 1pt

\newfontfamily{\thaifont}[Scale=1.23]{TH Sarabun New}
\newfontfamily{\thaifontsf}{TH Sarabun New}
\newfontfamily{\thaifonttt}{TH Sarabun New}

%%% PAGE DIMENSIONS
\usepackage{geometry} % to change the page dimensions
\geometry{a4paper} % or letterpaper (US) or a5paper or....

\usepackage{graphicx} % support the \includegraphics command and options
% \usepackage{fonts-tlwg}

\usepackage[parfill]{parskip} % Activate to begin paragraphs with an empty line rather than an indent

%%% PACKAGES
\usepackage{booktabs} % for much better looking tables
\usepackage{array} % for better arrays (eg matrices) in maths
\usepackage{enumitem}
\usepackage{subfigure} % make it possible to include more than one captioned figure/table in a single float

%\newcommand*\chem[1]{\textup{#1}}

%\usepackage{hyperref}
%\hypersetup{
%  colorlinks,
%  citecolor=black,
%  filecolor=black,
%  linkcolor=blue,
%  urlcolor=black
%}

%% Chapter Heading %%%%%
%\usepackage{titlesec} %% tufte-book already loads titlesec
%% Drawing pics with tikz %%%
\usepackage{tikz}
\usetikzlibrary{arrows,calc,decorations,shapes,decorations.pathmorphing,patterns}
\usepackage{pgfplots}

\definecolor{lightblue}{RGB}{180,220,255}
\definecolor{titlepagecolor}{cmyk}{1,.60,0,.40}
\definecolor{namecolor}{cmyk}{1,.50,0,.10}

\titleformat{\chapter}
{\huge\bfseries}
{\thechapter}{20pt}{\huge}

%% Section Heading %%%%%

\titleformat{\section}
{\Large\bfseries}
{\thesection}{20pt}{}

% Table of contents font setting
\usepackage[subfigure]{tocloft}
\renewcommand{\cftpartfont}{\Large\bfseries}
\renewcommand{\cftchapfont}{\large\bfseries}
\renewcommand{\cftsecfont}{\bfseries}


\begin{document}

\frontmatter

\begin{titlepage}
  \newgeometry{top=1cm,left=1cm} %defines the geometry for the titlepage
  \pagecolor{SkyBlue}
  \raggedright \includegraphics[scale=0.2]{pictures/logo-tu} \\
  \noindent
  \color{white}
  \makebox[0pt][l]{\rule{1.3\textwidth}{1pt}}
  \par
  \noindent
  \textbf{มหาวิทยาลัยธรรมศาสตร์} \textcolor{DarkBlue}{คณะวิศวกรรมศาสตร์}
  \vfill
  \hspace{1cm}
  \vfill
  \noindent
  \color{black}
  \raggedleft{\Huge\textbf{พลังงานงานทดแทนในประเทศไทย}}
  \vskip\baselineskip
  \noindent
  {\huge\color{Black}{สัปปินันทน์ เอกอำพน}}
\end{titlepage}

\restoregeometry % restores the geometry
\pagecolor{White}% Use this to restore the color pages to white


\chapter*{คำนำ}

ตำราเล่มนี้ถูกเขียนขึ้นเพื่อใช้ประกอบการเรียนการสอนเกี่ยวกับการใช้พลังงานแสงอาทิตย์สำหรับนักศึกษาปี ๓ - ๔ และสำหรับบุคคลทั่วไปที่มีความสนใจทางด้านดังกล่าว โดยที่แม้เนื้อหาบางส่วนจะมีคณิตศาสตร์ชั้นสูงเพื่อช่วยในการแสดงความสัมพันธ์ระหว่างตัวแปร แต่ความตั้งใจหลักของผู้เขียนต้องการจะให้ผู้ที่มีความสนใจและมีพื้นฐานคณิตศาสตร์ระดับมัธยมปลายควรจะสามารถอ่านแล้วเข้าใจได้ ทั้งนี้เนื่องจากผู้เขียนเล็งเห็นความสำคัญของการสร้างความเข้าใจพื้นฐานเรื่องของพลังงานแสงอาทิตย์ รวมถึงเทคโนโลยีต่างๆที่จะนำไประยุกต์ใช้เพื่อกักเก็บ แปลง หรือนำพลังงานนี้ไปใช้ เพื่อให้ผู้อ่านจะได้มีความเข้าใจที่ถูกต้อง มีพื้นฐานความรู้ที่เหมาะสมในการทำงานในเทคโนโลยีพลังงานสะอาดในอนาคต หรือแม้แต่สามารถทำความเข้าใจและคำนึงถึงความเหมาะสมของนโยบายหรือโครงการที่เกี่ยวกับพลังงานแสงอาทิตย์ได้โดยไม่เชื่อเพียงคำโฆษณาหรืออวดอ้างที่อาจจะเกินความเป็นจริงในหลายครั้ง

ผู้เขียนหวังว่าข้อมูลที่ได้รับการรวบรวมไว้ในตำราเล่มนี้จะเป็นประโยชน์ต่อผู้อ่านในวงกว้าง มิใช่เฉพาะระดับนักศึกษาหรือนักวิชาการเท่านั้น อย่างไรก็ดี ถ้าหากผู้อ่านมีความรู้พื้นฐานทางด้านฟิสิกส์พื้นฐาน จะทำให้สามารถเข้าใจเนื้อหาและบทวิเคราะห์ได้ดียิ่งขึ้น รวมถึงสามารถนำความรู้ที่ได้รับนำไปวิเคราะห์ข้อมูลอื่นๆได้ด้วยตนเอง

\vspace{2cm}\hspace{5cm}สัปปินันทน์ เอกอำพน

\clearpage
\tableofcontents

\clearpage
\listoffigures

\clearpage
\listoftables

\mainmatter

\part{พลังงานทดแทนเบื้องต้น}

\chapter{พลังงานแสงอาทิตย์}

เวลาพูดถึงพลังงานแสงอาทิตย์นั้น หลายๆคนอาจจะนึกถึงแดดร้อนๆในช่วงเดือนมีนาคมหรือเมษายน แต่จริงๆแล้วจะรู้ไหมว่าพลังงานที่มีอยู่ในแสงอาทิตย์นั้นประกอบด้วยหลายส่วน การจะตักตวงพลังงานแสงอาทิตย์มาใช้ให้ได้เต็มที่นั้น จำเป็นที่เราจะต้องมีความเข้าใจถึงส่วนประกอบเหล่านี้

เนื่องจากพลังงานแสงอาทิตย์นั้นเป็นพลังงานที่เกิดขึ้นมาจากการแผ่รังสีของดวงอาทิตย์ออกมาในรูปของคลื่นแม่เหล็กไฟฟ้าในช่วงคลื่นต่างๆ ดังนั้นเราควรจะเริ่มทำความเข้าใจกับการแผ่รังสีของวัตถุดำก่อน

\section{การแผ่รังสีของวัตถุดำ (Blackbody Radiation)}

การแผ่รังสีของวัตถุดำเกิดจากการแผ่รังสีคลื่นแม่เหล็กไฟฟ้าจากความร้อนของวัตถุซึ่งอยู่ในสภาวะสมดุลทางอุณหพลศาสตร์กับสิ่งแวดล้อม ซึ่งช่วงความถี่และความเข้มข้นของคลื่นต่างๆนั้นขึ้นอยู่กับอุณหภูมิของวัตถุดังกล่าว อย่างไรก็ดี ในความเป็นจริงแล้วไม่มีวัตถุได้ที่มีการแผ่รังสีเหมือนวัตถุดำแท้จริง โดนเฉพาะอย่างยิ่งดาวฤกษ์อย่างพระอาทิตย์นั้นก็ไม่ได้อยู่ในสภาวะสมดุลกับสิ่งแวดล้อม แต่ความเข้าใจเรื่องของการแผ่รังสีนี้ก็สามารถนำมาใช้ทำความเข้าใจส่วนประกอบของแสงอาทิตย์ได้

ยกตัวอย่างเช่น ในวัตถุที่มีอุณหภูมิต่ำนั้น ในห้องมืดจะมองเห็นเป็นสีดำเนื่องจากช่วงคลื่นที่แผ่ออกมาเป็นช่วงอินฟราเรดซึ่งมองด้วยตาเปล่าไม่เห็น เมื่ออุณหภูมิสูงขึ้นถึงราว 500$^{\circ}$ C การแผ่รังสีเริ่มเข้าอยู่ในช่วงความถี่ที่ตามองเห็น (visible spectrum) และจะเริ่มมีสีแดง เมื่ออุณหภูมิสูงมากจะออกเป็นสีฟ้าขาว เมื่อวัตถุมีการแผ่รังสีเป็นสีขาว แสดงว่ามีการแผ่รังสีบางส่วนออกมาเป็นรังสีอัลตราไวโอเลต

ดวงอาทิตย์ซึ่งมีอุณหภูมิที่ผิวประมาณ 5800 K นั้น มีการแผ่รังสีออกมามากที่สุดในช่วงคลื่นแสงและอินฟราเรด และมีจำนวนอีกเล็กน้อยในช่วงอัลตราไวโอเลต

\section{ทิศทางของแสงอาทิตย์}

\chapter{เทคโนโลยีการติดตามดวงอาทิตย์}

เนื่องจากดวงอาทิตย์เคลื่อนที่อยู่ตลอดเวลา และพลังงานของแสงอาทิตย์ที่ตกกระทบลงบนพื้นที่หนึ่งๆขึ้นอยู่กับความเข้มข้นของแสงและมุมตกกระทบ เพื่อจะเพิ่มพลังงานแสงอาทิตย์ที่ได้รับ เราสามารถออกแบบอุปกรณ์ให้มีความสามารถในการติดตามดวงอาทิตย์ (solar tracking) ซึ่งในปัจจุบันมีเทคโนโลยีหลายวิธีที่ใช้ในการติดตาม ซึ่งแบ่งได้เป็น 2 ประเภทใหญ่

\section{การติดตามแบบใช้พลังงาน}

การติดตามดวงอาทิตย์แบบใช้พลังงานหรือที่เรียกว่า Active Tracking นั้นเป็นการใช้ระบบ Feedback Loop โดยใช้ตัวรับแสงเพื่อช่วยในการบอกตำแหน่งของดวงอาทิตย์ที่ประเมินผล แล้วส่งสัญญาณให้กับระบบควบคุมให้เคลื่อนที่ไปยังตำแหน่งที่ต้องการ

\section{การติดตามแบบไม่ใช้พลังงาน}

\chapter{เซลล์แสงอาทิตย์}

เซลล์แสงอาทิตย์ (Solar cell หรือ Photovoltaic cell) เป็นอุปกรณ์ที่สามารถแปลงพลังงานจากคลื่นแม่เหล็กไฟฟ้าในแสงอาทิตย์ให้เป็นพลังงานไฟฟ้าได้โดยตรงโดยใช้ปรากฏการณ์โฟโตโวลตาอิก (Photovoltaic effect) ปรากฏการณ์นี้เกิดขึ้นจาก

\chapter{พลังงานความร้อนแสงอาทิตย์}

\section{การรับความร้อนโดยตรง}

\section{การรับความร้อนแบบรวมแสง}

\subsection{การใช้เลนส์}

\subsection{การใช้กระจกรวมแสง}

\section{คุณสมบัติของตัวรับแสง}

\section{เทคโนโลยีนำความร้อนไปใช้ต่อ}

\subsection{เทอร์โมอิเลกตริก}

เทอร์โมอิเลกตริกซิตี้เป็นปรากฏการณ์ที่เกิดศักย์ไฟฟ้าขึ้นบนตัวนำหรือสารกึ่งตัวนำที่มีอุณหภูมิเปลี่ยนไป 

\subsection{เครื่องยนต์สเตอร์ลิง}

\part{การกักเก็บพลังงานแสงอาทิตย์}

\chapter{เทคโนโลยีการกักเก็บพลังงาน}

ปัญหาหนึ่งของพลังงานแสงอาทิตย์และพลังงานทดแทนเช่นพลังงานลมหรือพลังงานคลื่นคือความไม่แน่นอนและไม่สามารถควบคุมได้ ซึ่งเป็นเกณฑ์วัดสำคัญของการสามารถพึ่งพาแหล่งพลังงานชนิดหนึ่งๆได้ ยกตัวอย่างเช่น ในกรณีของโรงไฟฟ้าพลังงานแก็สธรรมชาติ สามารถเปิดต่อเนื่องตลอดเวลาได้และสามารถเพิ่มหรือลดกำลังการผลิดได้ตามอุปสงค์อย่างไม่ยากเย็นนัก ในทางตรงกันข้าม พลังงานแสงอาทิตย์ไม่สามารถผลิตต่อเนื่องตลอดเวลาได้เนื่องจากช่วงเลากลางวันและกลางคืน นอกจากนี้ยังมีเรื่องของเมฆ ความชื้นในอากาศ ดังนั้น หากต้องการจะสร้างโรงไฟฟ้าพลังงานแสงอาทิตย์ (ไม่ว่าจะเป็นแบบ photovoltaics หรือ solar thermal หรือ แบบอื่นๆ) จำเป็นจะต้องสร้างเผื่อความไม่แน่นอนเหล่านี้ เช่นถ้ามีความต้องการพลังงานไฟฟ้า 10 MW อาจจะต้องสร้างโรงไฟฟ้าที่สามารถผลิตได้ 20 MW แล้วมีการกักเก็บส่วนที่เกินความต้องการไว้ใช้ในยามที่ไม่มีแสงอาทิตย์หรือพลังงานไม่เพียงพอต่อความต้องการผลิตต่อความต้องการผลิตต่อความต้องการผลิตต่อความต้องการผลิตต่อความต้องการผลิตต่อความต้องการผลิตต่อความต้องการผลิตต่อความต้องการผลิตต่อความต้องการผลิต 

ดังนั้น การจะลดผลกระทบจากความผันผวนและเพิ่มความสามารถในการควบคุมแหล่งพลังงาน โดยเฉพาะอย่างยิ่งแหล่งพลังงานทดแทนเช่นพลังงานแสงอาทิตย์ จำเป็นที่จะต้องมีอุปกรณ์กับเก็บพลังงานที่มีประสิทธิภาพเพื่อเก็บพลังงานส่วนเกินไว้ แล้วสามารถดึงพลังงานที่กเก็บไว้มาใช้ในช่วงที่มีความต้องการได้โดยไม่ต้องพึ่งพาแหล่งพลังงานโดยตรง

วิธีการกับเก็บพลังงานจากแสงอาทิตย์สามารถแบ่งได้เป็นหลายประเภท ซึ่งแต่ละประเภทก็มีจุดเด่นและจุดด้อยต่างกันไป พึงคำนึงไว้เสมอว่าไม่มีเทคโนโลยีใดที่จะเหมาะสมและดีกว่าเทคโนโลยีอื่นในทุกสถานการณ์ เราจึงควรทำความเข้าใจประเด็นต่างๆที่สำคัญเหล่านี้ไว้ เพื่อจะได้นำเทคโนโลยีเหล่านี้ไปประยุกต์ใช้ในสถานการณ์ต่างๆได้อย่างเหมาะสม

\section{บ่อกักเก็บพลังงานแสงอาทิตย์ (Solar Ponds)}

บ่อกักเก็บพลังงานแสงอาทิตย์ในที่นี้หมายถึงบ่อกับเก็บของเหลวซึ่งสามารถกับเก็บความร้อนจากพลังงานแสงอาทิตย์เพื่อนำไปใช้ประโยชน์ต่อไปได้ ในปัจจุบันบ่อกักเก็บพลังงานแสงอาทิตย์ส่วนมากใช้สารละลายเกลือคลอไรด์หรือซัลเฟตในน้ำ หลักการทำงานของบ่อดังกล่าวคือการแบ่งชั้นของสารละลายตามความความเข้มข้น โดยสารละลายที่มีความเข้มข้นมากจะตกอยู่ที่ชั้นล่างเนื่องจากมีความหนาแน่นสูง และสารละลายที่มีความเข้มข้นน้อยจะลอยอยู่ด้านบนเนื่องจากมีความหนาแน่นน้อย ซึ่งการแบ่งชั้นนี้จะป้องกันการหมุนเวียนของสารละลายเมื่อได้รับความร้อน ซึ่งในบ่อน้ำปกติเมื่อได้รับความร้อน จะมีการหมุนเวียนขึ้นเนื่องจากน้ำที่ร้อนกว่าจะมีการขยายตัว ทำให้ความหนาแน่นลดลงและลอยขึ้นสู่ด้านบน แต่ในบ่อน้ำที่มีการแบ่งชั้นของสานละลายนี้จะไม่มีการหมุนเวียนของสารละลาย ทำให้สามารถกักเก็บความร้อนไว้ได้

กระบวนการสร้างบ่อกักเก็บพลังงานแสงอาทิตย์มีอยู่ 2 วิธี

\subsection{บ่อกับเก็บแบบประดิษฐ์}

บ่อกักเก็บพลังงานแบบนี้สร้างโดยการเติมสารละลายที่มีความเข้มข้นจากสูงลงไปสู่ชั้นล่างแล้วลดลงต่ำลงเมื่อเพิ่มระดับน้ำขึ้นเรื่อยๆ จนเมื่อเติมเสร็จ บ่อก็จะสามารถกับเก็บพลังงานแสงอาทิตย์ไว้ได้

\subsection{บ่อกักเก็บแบบเกิดเอง}

บ่อประเภทนี้อาศัยหลักการของการละลายอิ่มตัวของเกลือในน้ำที่อุณหภูมิต่างๆกัน โดยที่ความสามารถในการละลายแปรผันตรงกับอุณหภูมิของตัวทำละลาย ซึ่งเกลือที่จะนำมาใช้ในบ่อประเภทนี้ จำเป็นจะต้องมีอัตราการเปลี่่ยนแปลงความสามารถในการละลายต่ออุณหภูมิสูง เพื่อที่จะได้สามารถสร้าง gradient ของความเค็มต่อความลึกได้สูง และมีความสามารถในการเก็บความร้อนได้ดี
\section{แบตเตอรี่}

\section{ล้อตุนกำลัง (Flywheel)}

\chapter{พลังงานลม}

\chapter{พลังงานชีวภาพ}

\part{ศักยภาพของพลังงานทดแทนในประเทศไทย}

\chapter{การประเมินต้นทุนทางเศรษฐศาสตร์}

เคยสงสัยกันบ้างไหมว่า เวลาที่การไฟฟ้าเก็บค่าไฟเราหน่วยละ 3 บาทกว่าๆนั้น เขาคิดคำนวณกันมาอย่างไร มีหลักฐานอ้างอิงหรือข้อมูลอะไรมาช่วยสนับสนุนนี้ไหม หรือว่าแค่นั่งเทียนกำหนดเลขกลมๆขึ้นมา จริงๆแล้วก็คงไม่ใช่อย่างนั้น และแน่นอนว่าค่าไฟที่เก็บนั้นก็คงไม่ได้เท่าทุนพอดี คงจะต้องมีส่วนบวกเพื่อให้เป็นกำไรไว้ไม่มากก็น้อยเป็นแน่

ในบทนี้ เราจะมาพูดถึงการวิเคราะห์ต้นทุนทางเศรษฐศาสตร์ของการผลิตพลังงาน เพื่อใช้ในการเปรียบเทียบระหว่างการลิตไฟฟ้าปัจจุบัน (พ.ศ. 2560) ซึ่งโดยส่วนมากยังพึ่งพาเชื้อเพลิงปิโตรเลียมอยู่่ กับการผลิตไฟฟ้าจากพลังงานทดแทนซึ่งเราได้กล่าวถึงเทคโนโลโยีและอุปกรณ์ที่ต้องใช้ไปในส่วนที่ 1

หลายครั้งที่วิศวกรโดยเฉพาะอย่างยิ่งในสถานศึกษา (ตัวผมเองก็ด้วย) คิดวิเคราะห์ปัญหาทางพลังงานที่มีอยู่ในปัจจุบันโดยยังไม่ได้พิจารณาเรื่องของความเหมาะสมของเทคโนโลยีทางเศรษฐศาสตร์ หรือที่เรียกง่ายๆว่า เทคโนโลยีนั้นมันแพงเกินไปหรือเปล่า การจะพิจารณาความเป็นไปได้ที่จะนำเทคโนโลยีพลังงานหนึ่งๆมาใช้ แม้ว่าจะมีความล้ำสมัย สะอาด และเป็นมิตรต่อสิ่งแวดล้อมเพียงใด หากมีราคาแพงกว่าของเดิมที่ใช้อยู่ปัจจุบัน ก็คงยากที่จะโน้มน้าวให้ประชาชนส่วนมากเห็นดีเห็นงามไปด้วย ไม่ใช่ว่าพวกเขาไม่ได้รักโลก หรือไม่ห่วงเรื่องสิ่งแวดล้อม แต่ว่าการจะบอกว่าได้โปรดใช้ของที่แพงขึ้นหน่อยเพื่อให้โลกสะอาดขึ้นก็ฟังดูเป็นข้ออ้างที่อาจจะดูหลักลอยไปสักหน่อย วิธีง่ายที่สุดที่จะชวนให้ประชาชนทั่วไปหันมาสนใจการใช้พลังงานทดแทนอย่างจริงจังก็คือต้องบอกว่าของใหม่นั้น\textbf{ถูกกว่า}

ดังนั้น เพื่อจะแน่ใจว่าเทคโนโลยีพลังงานทดแทนของเรานั้นถูกกว่าไฟฟ้าที่ผลิตอยู่ปัจจุบัน เราจำเป็นจะต้องทำความเข้าใจก่อนว่าโครงสร้างต้นทุนการผลิตไฟฟ้า หรือพลังงานอื่นๆที่ใช้ในครัวเรือนปัจจุบันนั้นเป็นอย่างไร

\section{โครงสร้างต้นทุน}

ศาสตร์เรื่องการวิเคราะห์โครงสร้างต้นทุนนั้นมีมานานโขอยู่ เริ่มจากปี ... ซึ่งพลังงานก็นับเป็นผลิตภัณฑ์อย่างหนึ่งซึ่งใช้สามารถจะวิเคราะห์ต้นทุนได้ การแบ่งประเภทต้นทุนนั้นสามารถทำได้อยู่หลายวิธี แล้วแต่จุดประสงค์และการนำไปใช้ประโยชน์ อย่างไรก็ดี ในหนังสือเล่มนี้เราต้องการศึกษาประเภทของต้นทุนเพื่อทำความเข้าใจแนวโน้มการเปลี่ยนแปลงเมื่อมีการพัฒนาเทคโนโลยีต่างๆที่เปลี่ยนไป จึงได้เลือกใช้วิธีการจำแนกต้นทุนตามความสัมพันธ์กับระดับของกิจกรรม ซึ่งสามารถสะท้อนความเปลี่ยนแปลงอันขึ้นอยู่กับระดับการผลิต โดยโครงสร้างต้นทุนแบบนี้สามารถแบ่งออกเป็นประเภทดังนี้

\begin{enumerate}
\item ต้นทุนคงที่ (Fixed Costs)

  เป็นต้นทุนส่วนที่ไม่มีการเปลี่ยนแปลงในช่วงระดับการผลิตหนึ่ง ซึ่งทำให้ต้นทุนต่อหน่วยลดลงเมื่อเพิ่มปริมาณการผลิตมากขึ้น
\item ต้นทุนผันแปร (Variable Costs)

  เป็นต้นทุนส่วนที่ต้นทุนรวมมีการเปลี่ยนแปลงขึ้นอยู่กับปริมาณการผลิต ในขณะที่ต้นทุนต่อหน่วยยังคงที่

\item ต้นทุนผสม (Mixed Costs)

  เป็นต้นทุนที่มีลักษณะของทั้งต้นทุนคงที่และผันแปรผสมกัน สามารถแบ่งได้เป็นสองประเภท

  \begin{enumerate}
  \item ต้นทุนกึ่งผันแปร (semi variable cost) เป็นต้นทุนที่จะมีส่วนหนึ่งคงที่ทุกระดับกิจกรรม และมีส่วนที่ผันแปรไปกับระดับกิจกรรม เช่น ค่าโทรศัพท์ เป็นต้น บางครั้งก็เป็นการยากที่จะประเมินส่วนที่คงที่หรือแปรผันของส่วนนี้
  \item ต้นทุนเชิงขั้น (step cost) หรือต้นทุนกึ่งคงที่ (semi fixed cost) หมายถึงต้นทุนที่คงที่ในช่วงระดับกิจกรรมหนึ่ง และเปลี่ยนไปคงที่ในอีกระดับกิจกรรมหนึ่ง เช่น ค่าผู้ควบคุมงาน เงินเดือน
  \end{enumerate}
\end{enumerate}

\chapter{การพัฒนาอย่างยั่งยืน}

\backmatter

\end{document}
%%% Local Variables: 
%%% coding: utf-8
%%% mode: latex
%%% TeX-engine: xetex
%%% End: 
